% ============================================================
%  Polinômios — Apostila Premium | PratiqueJá
%  Engine: XeLaTeX
% ============================================================
\documentclass[12pt]{article}
\usepackage{../../pratiqueja}

\newcommand{\hlm}[1]{{\color{darkblue}#1}}

\setsubject{Polinômios}

\begin{document}
\pagenumbering{arabic}
\setstretch{1.2}
\color{bodytext}

\heroheader{Polinômios}%
           {Operações, fatoração e Teorema do Resto}%
           {Matemática · Avançado}

\setlength{\columnsep}{18pt}
\setlength{\columnseprule}{0.5pt}
\def\columnseprulecolor{\color{exborder}}

\begin{multicols}{2}

\TW{Def}{darkblue}{Definição e Nomenclatura}{Grau, coeficientes e termos semelhantes}

\regra{
Um \hl{polinômio} em $x$ é uma expressão da forma
\[
  p(x) = a_n x^n + a_{n-1} x^{n-1} + \cdots + a_1 x + a_0
\]
O \hl{grau} é o maior expoente com coeficiente não nulo.
}

\begin{itemize}
  \item \textbf{Monômio} — 1 termo \quad ($3x^2$)
  \item \textbf{Binômio} — 2 termos \quad ($x^3 - 5$)
  \item \textbf{Trinômio} — 3 termos \quad ($x^2 + 4x - 1$)
  \item \textbf{Polinômio} — 4 ou mais termos
\end{itemize}

\textbf{Termos semelhantes}: mesma parte literal — podem ser somados: $3x^2 + 5x^2 = 8x^2$.

\exemp{
$p(x) = 3x^4 - 2x^2 + 5x - 1$: \hl{grau 4}, coef. líder $a_4 = 3$, termo ind. $a_0 = -1$

$4x^3 - x^2 + 3x^3 + 5x^2 - 2x = \hlm{7x^3 + 4x^2 - 2x}$ (grau 3)
}

\nota{O \textbf{polinômio nulo} (todos coef. = 0) não tem grau definido. O polinômio constante não nulo tem grau 0.}

\TW{Op}{darkblue}{Operações com Polinômios}{Adição, subtração e multiplicação}

\regra{
\textbf{Adição/Subtração:} agrupar termos semelhantes.

\textbf{Multiplicação:} lei distributiva — somar expoentes de mesma base.
}

\exemp{
$p = 3x^3 - x^2 + 4x - 2$ \quad $q = -x^3 + 2x^2 - x + 5$

$p+q = 2x^3 + x^2 + 3x + 3$

$p-q = 4x^3 - 3x^2 + 5x - 7$

$(x+2)(x^2-3x+1) = x^3-3x^2+x+2x^2-6x+2$

$\hlm{= x^3 - x^2 - 5x + 2}$

$(2x-1)(3x^2+x-4) = 6x^3+2x^2-8x-3x^2-x+4$

$\hlm{= 6x^3 - x^2 - 9x + 4}$
}

\TW{CN}{badgepurple}{Casos Notáveis}{Produtos especiais com resultado memorável}

\regra{
$(a + b)^2 = a^2 + 2ab + b^2$

$(a - b)^2 = a^2 - 2ab + b^2$

$(a + b)(a - b) = a^2 - b^2$
}

\exemp{
$(2x+3)^2 = 4x^2 + 12x + 9$

$(3x-1)^2 = 9x^2 - 6x + 1$

$(x+5)(x-5) = x^2 - 25$

$(2x+7)(2x-7) = 4x^2 - 49$
}

\nota{\textbf{Atalho ENEM:} $51 \times 49 = (50+1)(50-1) = 2500 - 1 = 2499$.}

\TW{Fat}{badgepurple}{Fatoração de Polinômios}{Fator comum, agrupamento, trinômio perfeito e diferença de quadrados}

\regra{
Fatorar é escrever o polinômio como produto de fatores de grau menor. As quatro técnicas aplicadas em sequência:

\textbf{1. Fator comum em evidência:} MDC dos coef. e menor potência da variável.

\textbf{2. Agrupamento:} agrupar pares que revelam fator comum.

\textbf{3. Trinômio perfeito:} $a^2 \pm 2ab + b^2 = (a \pm b)^2$.

\textbf{4. Diferença de quadrados:} $a^2 - b^2 = (a+b)(a-b)$.
}

\exemp{
\textbf{1. Fator comum:}

$6x^3 - 4x^2 + 2x = \hlm{2x(3x^2 - 2x + 1)}$

$x^3 - x = x(x^2-1) = \hlm{x(x+1)(x-1)}$

\textbf{2. Agrupamento:}

$x^3+2x^2+x+2 = x^2(x+2)+1(x+2) = \hlm{(x^2+1)(x+2)}$

\textbf{3. Trinômio perfeito:}

$x^2+6x+9 = \hlm{(x+3)^2}$

$4x^2-20x+25 = \hlm{(2x-5)^2}$

\textbf{4. Diferença de quadrados:}

$4x^2-9 = \hlm{(2x+3)(2x-3)}$

$x^4-16 = (x^2+4)(x^2-4) = \hlm{(x^2+4)(x+2)(x-2)}$
}

\TW{BR}{darkblue}{Divisão de Polinômios — Briot-Ruffini}{Algoritmo de divisão por $(x - a)$}

\regra{
O algoritmo de \hl{Briot-Ruffini} divide $p(x)$ pelo binômio $(x-a)$ usando apenas os coeficientes. Resultado: quociente $q(x)$ (grau $n-1$) e resto $r$.

\textbf{Passo a passo:} escreva os coef. em ordem decrescente (0 para potências ausentes); anote $a$ à esquerda; desça o 1º coef.; multiplique por $a$ e some ao próximo coef.; repita. O último valor é o resto.
}

\exemp{
\textbf{Divida $p(x) = x^3 - 6x^2 + 11x - 6$ por $(x-1)$:}

\vspace{4pt}
\noindent
\begin{minipage}[t]{\linewidth}
\footnotesize
\begin{tabular}{r|rrrr}
$1$ & $1$ & $-6$ & $11$ & $-6$ \\
  &   & $1$ & $-5$ & $6$ \\
\hline
  & $1$ & $-5$ & $6$ & $\mathbf{0}$ \\
\end{tabular}
\end{minipage}

\vspace{4pt}
Quociente: $x^2 - 5x + 6 = (x-2)(x-3)$

$\hlm{p(x) = (x-1)(x-2)(x-3)}$

\textbf{Divida $2x^3 - 3x^2 - 8x + 12$ por $(x-2)$:}

\vspace{4pt}
\noindent
\begin{minipage}[t]{\linewidth}
\footnotesize
\begin{tabular}{r|rrrr}
$2$ & $2$ & $-3$ & $-8$ & $12$ \\
  &   & $4$ & $2$ & $-12$ \\
\hline
  & $2$ & $1$ & $-6$ & $\mathbf{0}$ \\
\end{tabular}
\end{minipage}

\vspace{4pt}
$2x^2+x-6 = (2x-3)(x+2)$

$\hlm{p(x) = (x-2)(2x-3)(x+2)}$
}

\nota{\textbf{Potência ausente:} para $p(x) = x^3 - 4x + 1$, use coef. $1, 0, -4, 1$ na tabela.}

\TW{Teo}{darkblue}{Teorema do Resto e das Raízes}{$p(a) = 0 \;\Leftrightarrow\; (x-a)$ é fator de $p(x)$}

\regra{
Ao dividir $p(x)$ por $(x-a)$, o \hl{resto é igual a $p(a)$}:
\[
  p(x) = (x-a) \cdot q(x) + r \quad\Longrightarrow\quad r = p(a)
\]
\[
  p(a) = 0 \;\Longleftrightarrow\; (x-a) \text{ é fator de } p(x)
\]
}

\exemp{
\textbf{Resto de $q(x) = 2x^3+x^2-3x+4$ por $(x-1)$:}

$r = q(1) = 2+1-3+4 = \hlm{4}$

\textbf{Verifique que $x=1$, $x=-2$, $x=3$ são raízes de $x^3-2x^2-5x+6$:}

$p(1) = 1-2-5+6 = 0$ \;\ok \quad $(x-1)$ é fator

$p(-2) = -8-8+10+6 = 0$ \;\ok \quad $(x+2)$ é fator

$p(3) = 27-18-15+6 = 0$ \;\ok \quad $(x-3)$ é fator

$\hlm{p(x) = (x-1)(x+2)(x-3)}$

\textbf{Para qual $k$ é $x=2$ raiz de $x^3-kx+2$?}

$p(2) = 8 - 2k + 2 = 0 \Rightarrow 2k = 10$

$\hlm{k = 5}$
}

\nota{\textbf{Estratégia:} para fatorar um grau 3, procure raiz entre os divisores do termo independente, aplique Briot-Ruffini e fatore o quociente de grau 2 com Bhaskara se necessário.}

\end{multicols}

\end{document}
